\documentclass[../../../include/open-logic-section]{subfiles}

\begin{document}

\olfileid{sth}{spine}{foundation}

\olsection{આધાર}

આપણે માત્ર \emph{લગભગ} પૂર્ણ થયા છીએ, \emph{સંપૂર્ણપણે}
નહીં: $\ZFminus$માં એવું કશું નથી જે ખાતરી આપે કે
\emph{દરેક} ગણ કોઈ $V_\alpha$માં છે, એટલે કે દરેક ગણ કોઈક
તબક્કે રચાય છે.

ગણિતની દૃષ્ટિએ પદાનુક્રમની બહાર ગણો હોય કે નહીં તેની આપણે
\emph{ચિંતા કરવાની જરૂર નથી} એમ કહી શકાય. (હોય તો તેમને
અવગણી શકીએ.) પરંતુ દરેક ગણ કોઈક તબક્કે રચાય છે એવા વિચારથી
આપણે ગણની \emph{સંકલ્પના}ને પ્રેરિત કરી છે
(\olref[z][story]{sec}માં \stageshier{} જુઓ). તેથી
પદાનુક્રમની બહાર ગણો હોવાની શક્યતા દૂર કરવી છે. એ માટે
દરેક ગણ પદાનુક્રમમાં ક્યાંક આવે તેની ખાતરી આપતું નવું
સ્વયંસિદ્ધિ ઉમેરવું પડશે.

$V_\alpha$ આપણા તબક્કાઓ હોવાથી નીચેનું સીધું જ સ્વયંસિદ્ધિ
તરીકે ઉમેરી શકાય એમ વિચારીએ:

\begin{defish}
\emph{નિયમિતતા.} $\forall A \exists \alpha\, A \subseteq V_\alpha$
\end{defish}

આ રીત સંપૂર્ણપણે વાજબી હોત. પરંતુ આગળના વિભાગમાં સમજાવાનાં
કારણોસર આપણે તેને બદલે બીજું સ્વયંસિદ્ધિ લઈશું:

\begin{axiom}[આધાર]
$(\forall A \neq \emptyset)(\exists B \in A)A \cap B = \emptyset$.
\end{axiom}

થોડો પ્રયત્ન કરીને $\ZFminus$માં બતાવી શકાય કે આધારથી
નિયમિતતા મળે છે:
\begin{defn}
દરેક ગણ $A$ માટે લો:
\begin{align*}
	\text{cl}_0(A) &= A,\\
	\text{cl}_{n+1}(A) &= \bigcup \text{cl}_n(A),\\
	\text{trcl}(A) &= \bigcup_{n < \omega} \text{cl}_{n}(A).
\end{align*}
$\text{trcl}(A)$ને $A$નું \emph{પરંપરિત સંવરણ} કહીએ.
\end{defn}
\noindent
``પરંપરિત સંવરણ'' નામ યોગ્ય છે:
\begin{prop}\ollabel{subsetoftrcl}
$A \subseteq \trcl{A}$ અને $\trcl{A}$ પરંપરિત ગણ છે.
\end{prop}

\begin{proof}
સ્પષ્ટ છે કે $A = \text{cl}_0(A) \subseteq \text{trcl}(A)$.
વળી $x
\in b \in \trcl{A}$ હોય, તો કોઈ $n$ માટે
$b \in \text{cl}_n(A)$, તેથી $x
\in \text{cl}_{n+1}(A) \subseteq \text{trcl}(A)$.
\end{proof}

\begin{lem}\ollabel{lem:TransitiveWellFounded}
જો $A$ પરંપરિત ગણ હોય, તો કોઈ $\alpha$ માટે $A
\subseteq V_\alpha$.
\end{lem}

\begin{proof}
\olref[ordinals][opps]{defsupstrict}માંની
``$\supstrict(X)$''ની વ્યાખ્યા યાદ કરીને બે ગણો વ્યાખ્યાયિત કરો:
\begin{align*}
	D  &= \Setabs{x \in A}{\forall \delta\ x \nsubseteq V_\delta}\\
	\alpha &= \supstrict\Setabs{\delta}{(\exists x \in A)
	(x \subseteq V_\delta \land (\forall \gamma \in \delta)x \nsubseteq V_\gamma)}
\end{align*}
ધારો કે $D = \emptyset$. તેથી $x \in A$ હોય, તો એવું કોઈ
$\delta$ છે કે $x \subseteq V_\delta$ અને ક્રમવાચક સંખ્યાઓના
સુક્રમથી $(\forall \gamma \in \delta)x \nsubseteq V_\gamma$;
આથી $\delta \in \alpha$, અને
\olref[spine][Valphabasic]{Valphabasicprops} મુજબ
$x \in V_\alpha$. તેથી જોઈતું $A \subseteq
V_{\alpha}$ મળે છે.

તેથી $D = \emptyset$ બતાવવું પૂરતું છે. વિરોધાભાસ માટે ધારો કે
એવું નથી. આધાર મુજબ એવું $B \in D \subseteq A$ છે કે $D \cap B
= \emptyset$. જો $x \in B$ હોય, તો $A$ પરંપરિત હોવાથી
$x \in A$; અને $x \notin D$ હોવાથી
$\exists \delta\ x \subseteq
V_\delta$. તેથી હવે લો:
\[
	\beta = \supstrict\Setabs{\delta}{(\exists x \in B)(x \subseteq V_\delta \land (\forall \gamma < \delta)x \nsubseteq V_\gamma)}.
\]
\sourcecorrection{OLMD-119}{મૂળના આ સૂત્રમાં $b$ લખાયું છે,
પણ અહીં એવો કોઈ ચલ પસંદ કરાયો નથી; અગાઉ પસંદ કરેલો ગણ $B$
છે. તેથી મર્યાદિત અસ્તિત્વ-પરિમાણકમાં $b$ના સ્થાને $B$
મૂક્યું છે.}
અગાઉની જેમ $B \subseteq V_\beta$ મળે છે, જે $B \in
D$ સાથે વિરોધાભાસ છે.
\end{proof}

\begin{thm}\ollabel{zfentailsregularity}
નિયમિતતા સાચી છે.
\end{thm}

\begin{proof}
$A$ નિશ્ચિત લો. \olref{subsetoftrcl} મુજબ
$A \subseteq \trcl{A}$ અને આ સંવરણ પરંપરિત છે. તેથી
\olref{lem:TransitiveWellFounded} મુજબ કોઈ $\alpha$ માટે
$A \subseteq \trcl{A} \subseteq V_\alpha$.
\end{proof}

આ પરિણામો બતાવે છે કે $\ZFminus$માં શરતી વિધાન
$\emph{Foundation}\Rightarrow\emph{Regularity}$ સાબિત થાય છે.
\olref[spine][rank]{zfminusregularityfoundation}માં બતાવીશું કે
$\ZFminus$માં $\emph{Regularity}\Rightarrow\emph{Foundation}$
પણ સાબિત થાય છે. તેથી $\ZFminus$ના સંદર્ભમાં આધાર અને
નિયમિતતા \emph{સમકક્ષ} છે. એટલે $\ZFminus$ સ્વીકાર્યા પછી,
આધારનું ઔચિત્ય નિયમિતતા સાથેની સમકક્ષતાથી મળે છે; અને
નિયમિતતાનું ઔચિત્ય \stageshier{} પરથી સીધું મળે છે.

\end{document}
